\section{Feasibility}
The project is sized to finish within the semester on hardware we already have. All work runs on one workstation with two 48 GB GPUs, 120 CPU cores, and 754 GB of memory, with long jobs kept in detached terminal sessions. Nothing here trains a large model. The language model used for interpretation is an open-weight model run locally for inference only, so there are no API costs and no data leaves the machine.

The data is manageable. Datasets 1 and 3 are small enough to process in minutes. Dataset 2 is about 2.3 million cells, which is roughly 16 GB as a dense matrix of 2,000 selected genes in single precision and well under 100 GB in sparse form, so it fits in host memory and streams to the GPUs in batches. The costly step is rescoring candidate pairs, which we estimate at a few GPU hours for the full atlas, with the rest of the pipeline taking minutes to an hour. Our target is a complete run on Dataset 2 inside a single twelve-hour window, and we will confirm the estimate on a 100,000-cell subset before committing to the full run.

Work is divided so that only one member needs the workstation. Utshab writes the code and runs every experiment. The other members handle dataset documentation, checking proposed cluster names against marker databases through their web interfaces, building the literature comparison tables, and writing. Results are shared as tables and figures in the repository, so no one else has to run code.

The plan for the rest of the semester is as follows.
\begin{itemize}
\item \textbf{Late September.} Set up the software, download and preprocess the data, and reproduce the paper on Zeisel and PBMC14K.
\item \textbf{October.} Run the baselines on all three datasets, add the label-free selection rule, and run the diagnostic experiments on what the similarity measures.
\item \textbf{November.} Build and test the scalable neighbor graph, run the atlas-scale experiments, and set up the retrieval and annotation comparison.
\item \textbf{December.} Final experiments, ablations, writing, and the presentation.
\end{itemize}

We have fallbacks for the parts that could fail. If the Allen atlas download or its license proves awkward, the Cao 2020 human fetal atlas is a single file with a permissive license. If the candidate lists miss too many true neighbors, we increase the candidate count, which raises cost linearly rather than quadratically. If the network inference step is too slow to run for every cluster, we restrict it to one part of the label hierarchy or to aggregated cells, which is the approach the original authors use. If a full atlas run does not finish in time, the sketch-and-propagate result still answers RQ2 with a clearly stated approximation.
